\section{Planteamiento del Problema}

La documentaci\'on de hardware t\'ecnico complejo — sistemas electromec\'anicos con m\'ultiples componentes integrados, como el dron Holybro X500~V2 — sigue apoyada principalmente en im\'agenes 2D, planos PDF y l\'aminas de ensamblaje. Estos formatos est\'aticos presentan un l\'imite funcional claro: obligan al usuario a reconstruir mentalmente las relaciones espaciales entre piezas, lo que incrementa la demanda cognitiva durante tareas de inspecci\'on y consulta t\'ecnica (Sweller, 1988).

El paso hacia entornos 3D interactivos en el navegador no es inmediato. Los modelos CAD de origen tienen una geometr\'ia muy densa, incompatible con las limitaciones de memoria y rendimiento de una aplicaci\'on web. Convertirlos en un visor web operativo requiere un proceso de optimizaci\'on art\'iculado — retopolog\'ia, \textit{baking} de texturas, control del presupuesto de renderizado — que a\'un no est\'a documentado con criterios verificables para casos de hardware complejo en contexto acad\'emico.

El problema central puede resumirse en tres restricciones concretas:

\begin{itemize}
    \item \textbf{Tama\~no de los modelos.} Los archivos CAD originales pueden pesar varias decenas o cientos de megabytes; sin reducci\'on geom\'etrica y empaquetado eficiente, la carga en navegador resulta inviable.
    \item \textbf{Rendimiento en tiempo real.} Mantener una experiencia fluida (30 FPS o m\'as) exige controlar la cantidad de objetos, materiales y estados de renderizado activos simult\'aneamente.
    \item \textbf{Alternativas con limitaciones.} Las soluciones de \textit{streaming} remoto (como Unreal Pixel Streaming) ofrecen alta calidad visual pero dependen de servidores con GPU y presentan latencia, lo que las hace poco pr\'acticas para un prototipo acad\'emico de acceso directo desde el navegador.
\end{itemize}

Existe, por tanto, la necesidad de construir y evaluar un \textit{pipeline} de optimizaci\'on que permita desplegar hardware complejo en un visor 3D web, con capacidad de inspecci\'on t\'ecnica real, rendimiento controlado y accesibilidad desde cualquier navegador compatible.

\subsection{Pregunta de Investigaci\'on}

\noindent \textbf{Pregunta principal.} \textquestiondown Qu\'e diferencias se observan en el desempe\~no de tareas, la usabilidad percibida y la carga de trabajo cuando usuarios con perfil t\'ecnico-acad\'emico realizan actividades de exploraci\'on e identificaci\'on del dron Holybro X500~V2 mediante un prototipo web 3D interactivo, comparado con un soporte 2D de referencia, y bajo qu\'e condiciones t\'ecnicas ese prototipo resulta viable en el navegador?

\newpage
